\section{Operator: Aritmetika, Perbandingan, Logika, Penugasan, dan String}

Operator adalah simbol yang digunakan untuk melakukan operasi pada nilai dan variabel. PHP menyediakan berbagai jenis operator yang lengkap untuk kebutuhan pemrograman, mulai dari perhitungan matematika sederhana hingga operasi logika kompleks. Memahami cara kerja setiap operator dan urutan prioritasnya (\textit{operator precedence}) sangat penting untuk menulis ekspresi yang menghasilkan nilai yang benar \cite{phpManual}.

\textbf{Operator aritmetika} digunakan untuk operasi matematika dasar. PHP mendukung penjumlahan (\code{+}), pengurangan (\code{-}), perkalian (\code{*}), pembagian (\code{/}), modulo/sisa bagi (\code{\%}), dan pangkat (\code{**}). Operator pembagian \code{/} menghasilkan float jika hasilnya bukan bilangan bulat, sedangkan modulo selalu mengembalikan sisa pembagian. Operator pangkat \code{**} diperkenalkan di PHP 5.6 dan sangat berguna untuk perhitungan matematika \cite{w3schoolsPhp}.

\textbf{Operator perbandingan} digunakan untuk membandingkan dua nilai dan menghasilkan boolean (\code{true} atau \code{false}). Yang paling penting untuk diingat adalah perbedaan antara \code{==} (sama nilai, dengan type juggling) dan \code{===} (identik: sama nilai \textbf{dan} sama tipe). Menggunakan \code{===} adalah praktik yang lebih aman karena menghindari konversi tipe implisit yang dapat menyebabkan bug tersembunyi.

\begin{table}[ht]
\centering
\begin{tabular}{|P{1.5cm}|P{3.5cm}|P{3.5cm}|P{3.5cm}|}
\hline
\textbf{Kategori} & \textbf{Operator} & \textbf{Contoh} & \textbf{Hasil} \\
\hline
\multirow{6}{*}{Aritmetika} & \code{+} Penjumlahan & \code{5 + 3} & \code{8} \\
\cline{2-4}
& \code{-} Pengurangan & \code{10 - 4} & \code{6} \\
\cline{2-4}
& \code{*} Perkalian & \code{3 * 7} & \code{21} \\
\cline{2-4}
& \code{/} Pembagian & \code{10 / 3} & \code{3.333...} \\
\cline{2-4}
& \code{\%} Modulo & \code{10 \% 3} & \code{1} \\
\cline{2-4}
& \code{**} Pangkat & \code{2 ** 8} & \code{256} \\
\hline
\multirow{4}{*}{Perbandingan} & \code{==} Sama nilai & \code{"5" == 5} & \code{true} \\
\cline{2-4}
& \code{===} Identik & \code{"5" === 5} & \code{false} \\
\cline{2-4}
& \code{!=} atau \code{<>} & \code{5 != 3} & \code{true} \\
\cline{2-4}
& \code{<=>} Spaceship & \code{5 <=> 3} & \code{1} \\
\hline
\multirow{3}{*}{Logika} & \code{\&\&} / \code{and} & \code{true \&\& false} & \code{false} \\
\cline{2-4}
& \code{||} / \code{or} & \code{true || false} & \code{true} \\
\cline{2-4}
& \code{!} Not & \code{!true} & \code{false} \\
\hline
\multirow{3}{*}{Penugasan} & \code{=} Assign & \code{\$x = 10} & \code{10} \\
\cline{2-4}
& \code{+=}, \code{-=} & \code{\$x += 5} & \code{\$x = \$x + 5} \\
\cline{2-4}
& \code{.=} Concat & \code{\$s .= "!"} & \code{\$s = \$s . "!"} \\
\hline
\end{tabular}
\caption{Rangkuman operator PHP berdasarkan kategori}
\end{table}

\textbf{Operator logika} digunakan untuk mengombinasikan kondisi boolean, sangat penting dalam percabangan dan validasi. Operator \code{\&\&} (AND) mengembalikan \code{true} hanya jika kedua operand benar; \code{||} (OR) mengembalikan \code{true} jika salah satu operand benar; \code{!} (NOT) membalik nilai boolean. PHP juga menyediakan operator \textbf{null coalescing} \code{??} yang mengembalikan operand kiri jika ada dan tidak null, sangat berguna untuk memberikan nilai default pada variabel yang mungkin belum terdefinisi \cite{phpRightWay}.

\begin{webcode}[language=PHP, caption={Contoh penggunaan operator}]
<?php
// Aritmetika
$total = 100 + 50;        // 150
$sisa = 17 % 5;           // 2
$pangkat = 2 ** 10;       // 1024

// Perbandingan ketat vs longgar
var_dump(0 == "abc");     // true (type juggling!)
var_dump(0 === "abc");    // false (tipe berbeda)

// Null coalescing
$nama = $_GET["nama"] ?? "Anonim";  // default "Anonim"

// Spaceship operator (PHP 7+)
echo 1 <=> 2;  // -1 (kiri < kanan)
echo 2 <=> 2;  //  0 (sama)
echo 3 <=> 2;  //  1 (kiri > kanan)

// Concatenation (penggabungan string)
$sapaan = "Halo" . ", " . "dunia!";
echo $sapaan;  // Halo, dunia!
?>
\end{webcode}

\begin{konsep}
\textbf{Operator Spaceship (\code{<=>}):} Diperkenalkan di PHP 7, operator ini membandingkan dua nilai dan mengembalikan \code{-1}, \code{0}, atau \code{1}. Sangat berguna sebagai callback dalam fungsi pengurutan seperti \code{usort()}.
\end{konsep}

